\chapter{Introduction}

%====================================================
\section{1.1 Overview of the Chapter}

This chapter introduces the research background and foundational concepts related to explainable machine learning–based fault classification of power transformers using Dissolved Gas Analysis (DGA). The purpose of this chapter is to establish the context of the study, define the problem domain, and highlight the motivation for integrating machine learning and explainable artificial intelligence in transformer fault diagnosis.

Power transformers are essential components in modern electrical power systems, and their continuous operation is critical for ensuring reliable energy transmission and distribution. Due to their importance, even minor faults can lead to significant system instability and operational losses. Therefore, effective condition monitoring techniques are required to ensure early fault detection and system reliability \parencite{Wang2023Imbalanced}.

This chapter also provides a structured transition from traditional diagnostic approaches to data-driven machine learning methods. It explains how conventional DGA interpretation techniques are limited in handling complex fault patterns, and how machine learning models offer improved adaptability through data-driven learning mechanisms \parencite{Dladla2025SLR}.

Finally, the chapter outlines the key research gaps, objectives, and scope of the study, which collectively guide the development of the proposed methodology in subsequent chapters \parencite{Cao2024Advancement}.

%====================================================
\section{1.2 Background of Transformer Condition Monitoring}

Power transformers operate under continuous electrical, thermal, and mechanical stress throughout their operational lifespan. These stresses gradually degrade insulation materials, making transformers susceptible to internal faults if not properly monitored.

Transformer failures can lead to severe consequences such as power outages, equipment damage, and cascading failures across interconnected power networks. These failures may also impact critical infrastructure including healthcare systems, transportation networks, and industrial operations.

To mitigate such risks, condition monitoring systems are widely implemented in power systems. These systems aim to detect early signs of deterioration before catastrophic failure occurs. Among various diagnostic methods, Dissolved Gas Analysis (DGA) has been widely recognized as one of the most effective techniques for transformer health assessment \parencite{Ali2023ConventionalDGA}.

%====================================================
\section{1.3 Background of the Problem}

Dissolved Gas Analysis (DGA) is a diagnostic technique used to identify internal transformer faults by analyzing gases dissolved in transformer oil. When transformers experience thermal or electrical stress, decomposition of insulating materials produces characteristic gases that dissolve in the oil.

The key gases involved in DGA include hydrogen (H$_2$), methane (CH$_4$), ethane (C$_2$H$_6$), ethylene (C$_2$H$_4$), acetylene (C$_2$H$_2$), carbon monoxide (CO), and carbon dioxide (CO$_2$). These gases are directly associated with different fault mechanisms and energy levels inside transformers \parencite{IEC60599}.

International standards such as IEEE C57.104-2019 and IEC 60599:2022 define structured diagnostic guidelines for interpreting gas concentrations and ratios. These standards are widely used in industrial applications for transformer fault diagnosis \parencite{IEEEC57104}.

However, real-world transformer environments introduce complexity due to varying load conditions, aging effects, and environmental influences, which may affect gas generation patterns and reduce the effectiveness of rule-based interpretation methods.

%====================================================
\subsection{1.3.1 Physical Meaning of Gas Generation}

Gas generation in transformer oil is directly linked to the type and severity of internal faults. Different fault mechanisms produce distinct gas signatures depending on the energy level of the fault.

Low-energy electrical faults such as partial discharge typically generate hydrogen as the dominant gas due to localized ionization processes. Medium and high thermal faults generate hydrocarbon gases such as methane, ethane, and ethylene due to gradual decomposition of insulation materials.

High-energy electrical faults such as arcing generate acetylene, which is considered a critical indicator of severe internal transformer failure \parencite{Dladla2025SLR}.

Thermal faults are generally associated with gradual overheating processes, whereas electrical faults are associated with sudden high-energy discharge events. This physical relationship between gas generation and fault types forms the basis of DGA-based transformer diagnosis \parencite{Cao2024Advancement}.

%====================================================
\section{1.4 Conventional DGA Interpretation Methods}

Conventional DGA interpretation methods are rule-based diagnostic approaches used to classify transformer faults based on gas concentrations and ratios. These methods are widely used due to their simplicity and alignment with international standards.

Common methods include the Key Gas method, Rogers Ratio method, IEC Ratio method, Duval Triangle, and Duval Pentagon. These methods rely on predefined rules or graphical boundaries to classify transformer fault conditions \parencite{Ali2023ConventionalDGA}.

The main advantage of these methods is their ease of use and interpretability, making them suitable for field applications where quick decision-making is required.

However, these methods have limitations such as rigid classification boundaries, sensitivity to measurement errors, and inability to handle overlapping or complex fault conditions. In some cases, different methods may produce inconsistent results for the same transformer condition.

%====================================================
\section{1.5 Machine Learning Motivation in DGA Fault Diagnosis}

Machine learning techniques have been introduced to overcome the limitations of conventional DGA interpretation methods. These techniques enable data-driven learning of complex nonlinear relationships between gas concentrations and transformer fault types.

Unlike rule-based systems, machine learning models can automatically learn patterns from historical data without requiring predefined diagnostic rules. This allows them to adapt more effectively to real-world transformer operating conditions \parencite{Wang2021GPMC}.

One major advantage of machine learning is automation, which enables consistent and scalable fault classification across large datasets and monitoring systems.

However, challenges such as class imbalance, weak feature representation, and lack of interpretability remain significant issues in DGA-based machine learning systems \parencite{Azmi2025Imbalance}.

%====================================================
\subsection{1.5.1 Problem Formulation}

Despite the success of machine learning in transformer fault classification, several challenges remain unresolved.

Machine learning models often show inconsistent performance across different DGA datasets due to variations in data distribution and preprocessing methods. This reduces model reliability and generalization capability.

In addition, DGA datasets frequently suffer from class imbalance and weak feature representation, which negatively affects model learning performance for minority fault classes.

Furthermore, many machine learning models lack interpretability, making it difficult to understand how gas features influence final predictions \parencite{Laburu2024Insights}.

These challenges highlight the need for an integrated framework that combines feature optimization, consistent model evaluation, and explainable artificial intelligence techniques \parencite{Hu2025InterpretableEdge}.

%====================================================
\subsection{1.5.2 Research Gap}

The first research gap is the inconsistency in feature representation across studies, where different approaches use raw gas values, gas ratios, or hybrid features without standardized evaluation.

The second research gap is the lack of consistent benchmarking across machine learning models, making performance comparison unreliable across studies \parencite{Dladla2025SLR}.

The third research gap is the limited integration of explainable artificial intelligence techniques in DGA-based transformer fault classification, reducing interpretability in engineering applications \parencite{Zhang2022BiLevelXAI}.

%====================================================
\section{1.6 Hypotheses}

This study is guided by the following hypotheses:

\begin{itemize}
\item Different DGA feature sets significantly influence machine learning classification performance.
\item Traditional machine learning models show varying performance under identical datasets.
\item Explainable AI techniques can identify the most influential gas features contributing to transformer fault classification.
\end{itemize}

%====================================================
\section{1.7 Research Objectives}

The objectives of this study are:

\begin{itemize}
\item To evaluate machine learning models for DGA-based transformer fault classification.
\item To analyze the impact of different feature sets on classification performance.
\item To apply explainable artificial intelligence techniques to interpret model predictions.
\end{itemize}

%====================================================
\section{1.8 Scope of Work}

This study focuses on DGA-based transformer fault classification using benchmark datasets containing dissolved gas measurements.

The gases considered include hydrogen, methane, ethane, ethylene, acetylene, carbon monoxide, and carbon dioxide.

The study evaluates traditional machine learning models including SVM, Decision Tree, Random Forest, KNN, and boosting-based methods.

Deep learning approaches, real-time deployment systems, and hardware-level transformer monitoring are excluded from this study.

%====================================================
\section{1.9 Organisation of Thesis}

This thesis is organised into five chapters. Chapter 1 introduces the research background, problem formulation, and objectives. Chapter 2 presents a comprehensive literature review on DGA, machine learning, and explainable artificial intelligence. Chapter 3 describes the methodology including dataset preparation, feature engineering, model development, and evaluation methods. Chapter 4 presents results and discussion, while Chapter 5 concludes the study and provides future recommendations.